\section{An ultraviolet obstruction for the fixed-reference mean-field model}
\label{uv-section}

The cutoff dependence in Theorem~\ref{thm:mf} cannot be removed merely by
imposing a subcritical nuclear charge.  In the photon-free Hamiltonian fully
Wick ordered relative to the fixed free sea, a single unit point nucleus
already produces an energy gain proportional to minus the fermion cutoff.
The following construction uses an exactly neutral Slater determinant and
at most one coherent electron--positron excitation in addition to its
neutralizing electron.  Thus an unbounded number of pairs is not needed for
this obstruction.

Throughout this section $m>0$ and $0<\alpha<4/\pi$ are fixed,
$\Lambda:=\Lambda_f$, and the box side $L$ is fixed with
\begin{equation}
 L>\frac{\alpha c_*}{2m(1-\alpha\pi/4)},
 \qquad c_*=6+\frac{3\pi\sqrt6}{2}.
 \label{uv-box-range}
\end{equation}
The source is one point nucleus of charge $Z=1$ at the origin.  We retain the
same zero-mean Green function and nuclear self-energy convention as in the
main model.  There is therefore no nucleus--nucleus term in this example.
All assertions below concern the fixed bare coupling and fixed mass in
that explicitly specified Hamiltonian.

\begin{theorem}[No cutoff-independent mean-field relative bound for a unit point nucleus]
\label{uv-theorem}
Set $V=L^3$ and define the positive constants
\begin{align}
 b_0&=\frac1{21600\sqrt3\pi^2},\qquad
 c_1=c_*+\frac52,\qquad
 A_\alpha=\frac94+\frac{15\pi\alpha}{16},\nonumber\\
 c_\alpha&=\frac{A_\alpha\alpha^2b_0^2}
                    {A_\alpha^2+\alpha^2b_0^2}>0.
 \label{uv-constants}
\end{align}
For every cutoff satisfying
\begin{equation}
 \Lambda\ge\max\left\{
 32m,\ \frac{16\pi\sqrt3}{L},\
 4m+\frac{6\alpha c_1}{L},\
 \left(\frac{1024\pi\alpha}{9V}\right)^{1/3}
 \right\},
 \label{uv-cutoff-threshold}
\end{equation}
the Hartree--Fock energy, minimized over exactly neutral Slater determinants,
satisfies
\begin{equation}
 E^{\rm HF}_{\Lambda,L}(0,1)
       -E^{\rm HF,\vac}_{\Lambda,L}
 \le m-c_\alpha\Lambda.
 \label{uv-linear-collapse}
\end{equation}
Here $E^{\rm HF,\vac}_{\Lambda,L}=0$ by Theorem~\ref{thm:mf}.
Consequently no lower bound $-CK$, with $C$ independent of $\Lambda$, holds
for this fixed-reference mean-field model, even with $K=Z=1$; in particular,
this includes the physical coupling, for which $Z\alpha<2/\pi$.
The same obstruction applies to the relaxed mean-charge
variational class, which contains these Slater states.
\end{theorem}

The constants are deliberately conservative.  Their significance is their
strict positivity and independence of the cutoff, not their numerical size.
The proof below involves only finite-dimensional Fourier sums; no
unjustified interchange of ultraviolet and infinite-volume limits is needed.

\subsection{Opposite-energy Fourier packets with a positive Coulomb matrix element}
\label{uv-packets}

Let
\[
 \mathsf C=-i\mathsf a_1\mathsf a_2\mathsf a_3.
\]
The Dirac anticommutation relations imply that $\mathsf C$ is a Hermitian
unitary, commutes with all $\mathsf a_i$, and anticommutes with $\mathsf b$.
Choose a unit spinor $\xi$ satisfying
$\mathsf a_1\xi=\xi$ and $\mathsf C\xi=\xi$.  Such a spinor exists: the
commuting projections $(1+\mathsf a_1)/2$ and $(1+\mathsf C)/2$ have a
one-dimensional intersection, as follows by taking the trace of their
product.

Define a finite set of lattice momenta by
\begin{equation}
 \mathcal A_\Lambda=\left\{
 p\in\frac{2\pi}{L}\Z^3:
 \left|p_i-\frac\Lambda2\delta_{i1}\right|
        \le\frac\Lambda{8\sqrt3}\quad(i=1,2,3)
 \right\},\qquad N=|\mathcal A_\Lambda|.
 \label{uv-momentum-cube}
\end{equation}
Every such momentum obeys
\[
 \left|p-\frac\Lambda2e_1\right|\le\frac\Lambda8,
 \qquad \frac{3\Lambda}8\le|p|\le\frac{5\Lambda}8.
\]
The set and its reflection lie strictly inside the fermion cutoff and have
disjoint supports.  If $\Lambda L\ge16\pi\sqrt3$, elementary counting in
each of the three coordinate intervals gives
\begin{equation}
 N\ge\left(\frac{\Lambda L}{16\pi\sqrt3}\right)^3
       =\frac{V\Lambda^3}{12288\sqrt3\pi^3}.
 \label{uv-mode-count}
\end{equation}
Indeed an interval of length $\Lambda/(4\sqrt3)$ contains at least its
length divided by the lattice spacing, minus one, lattice points.  Under
the stated condition this is at least half that ratio.

For $p\in\mathcal A_\Lambda$ put
\[
 a_p=P_+(p)\xi,
 \qquad S_\Lambda=\sum_{p\in\mathcal A_\Lambda}\|a_p\|^2,
 \qquad \varphi_p(x)=V^{-1/2}e^{ip\cdot x},
\]
and define normalized spinors
\begin{equation}
 u(x)=S_\Lambda^{-1/2}\sum_{p\in\mathcal A_\Lambda}a_p\varphi_p(x),
 \qquad v(x)=\mathsf C u(-x).
 \label{uv-orbitals}
\end{equation}
To control the spinors quantitatively, note that
\[
 \left|\frac p{|p|}-e_1\right|\le\frac12,
 \quad
 1-\frac{|p|}{E(p)}\le\frac m{|p|},
 \quad \frac m{E(p)}\le\frac m{|p|}.
\]
Since $\mathsf a_1\xi=\xi$, these imply
\begin{equation}
 \|a_p-\xi\|
 \le\frac14+\frac m{|p|}
 \le\frac14+\frac{8m}{3\Lambda}\le\frac13
 \quad\text{if }\Lambda\ge32m.
 \label{uv-spin-bound}
\end{equation}
It follows that $4N/9\le S_\Lambda\le N$, so the definition of $u$ is
nondegenerate.

Let $\mathcal R$ be spatial inversion.  Since
$D_0\mathsf C\mathcal R=-\mathsf C\mathcal R D_0$, the two orbitals obey
\[
 P_+u=u,\qquad P_-v=v,\qquad
 \|u\|=\|v\|=1,\qquad \langle u,v\rangle=0.
\]
Moreover $G_L$ is even, so their diagonal Coulomb matrix elements are
exactly equal:
\begin{equation}
 \langle u,G_Lu\rangle=\langle v,G_Lv\rangle.
 \label{uv-equal-diagonal}
\end{equation}

The off-diagonal matrix element is positive and grows linearly.  For all
$p,r\in\mathcal A_\Lambda$, \eqref{uv-spin-bound} gives
\[
 \operatorname{Re}\langle a_p,\mathsf C a_r\rangle
 \ge1-\frac13-\frac13-\frac19=\frac29.
\]
No term $p+r$ vanishes.  The Fourier formula for $G_L$ therefore gives
\begin{align}
 b_\Lambda&:=\operatorname{Re}\langle u,G_Lv\rangle
   =\frac1{VS_\Lambda}\sum_{p,r\in\mathcal A_\Lambda}
       \frac{4\pi}{|p+r|^2}
       \operatorname{Re}\langle a_p,\mathsf C a_r\rangle\nonumber\\
 &\ge\frac{128\pi}{225}\frac{N}{V\Lambda^2}
 \ge b_0\Lambda.
 \label{uv-positive-matrix-element}
\end{align}
Here $|p+r|\le5\Lambda/4$, $S_\Lambda\le N$, and
\eqref{uv-mode-count} were used.  This establishes the required
electron--positron Coulomb coupling directly at finite volume.

\subsection{An exactly neutral Slater state}
\label{uv-neutral-state}

Write $P_u=|u\rangle\langle u|$ and similarly for other normalized orbitals.
Choose $0<t<\pi/2$, set $s=\sin t$, $c=\cos t$, and replace the occupied
negative-energy orbital $v$ by $z=cv+su$.  The resulting sea-relative
matrix is
\begin{equation}
 Q_t=s^2(P_u-P_v)+sc\bigl(|u\rangle\langle v|
                              +|v\rangle\langle u|\bigr),
 \qquad Q_t^2=s^2(P_u+P_v).
 \label{uv-pair-rotation}
\end{equation}
This rotation alone has exact charge zero.

To neutralize the unit nucleus, choose a unit spinor $\zeta$ with
$\mathsf b\zeta=\zeta$ and add the positive-energy zero-momentum orbital
$w(x)=V^{-1/2}\zeta$.  Its Fourier support is disjoint from those of $u,v$.
Then
\begin{equation}
 \gamma_t=P_- -P_v+P_z+P_w,
 \qquad Q^{\rm tot}_t=\gamma_t-P_-=Q_t+P_w
 \label{uv-full-projector}
\end{equation}
is an orthogonal projector of rank $2M_L+1$.  Its Slater determinant lies
in the exact neutral sector of the one-nucleus system.  In the free
electron--positron representation it is a superposition of a state with
one electron and a state with two electrons and one positron.

The kinetic contribution is
\begin{equation}
 \Tr(D_0Q^{\rm tot}_t)
 =m+2s^2\langle u,D_0u\rangle
 \le m+\left(\frac{5\Lambda}4+2m\right)s^2.
 \label{uv-kinetic}
\end{equation}
Equation~\eqref{uv-equal-diagonal} removes the order-$s^2$ part of the
pair's nuclear coupling.  The constant density of $w$ has zero coupling
to the zero-mean Green function.  Hence the complete electron--nucleus
expectation is
\begin{equation}
 \alpha D_L(\rho_{Q^{\rm tot}_t},\delta_0)
 =-2\alpha b_\Lambda sc
 \le-2\alpha b_0\Lambda sc.
 \label{uv-nuclear-attraction}
\end{equation}

\subsection{Bounds on every remaining Coulomb term}
\label{uv-coulomb}

The analytic Ewald certificate in the proof of Lemma~\ref{lem:pkato}
also supplies a convenient lower bound.  All real-space terms in
\eqref{eq:ewald-audit} are nonnegative, while
\eqref{eq:ewald-fourier-bound} bounds the absolute reciprocal-space sum
by $9/10$.  Therefore
\[
 G_1(x)>-\frac\pi9-\frac9{10}>-\frac54,
 \qquad |G_L(x)|\le G_L(x)+\frac5{2L}.
\]
The periodic Kato inequality thus implies, for either orbital $a=u,v$,
\begin{equation}
 \sup_y\int |G_L(x-y)|\,|a(x)|^2\dd x
 \le\frac\pi2\langle a,|p|a\rangle+\frac{c_1}{L}
 \le h_\Lambda,
 \qquad h_\Lambda=\frac{5\pi}{16}\Lambda+\frac{c_1}{L}.
 \label{uv-orbital-potential}
\end{equation}

Since $|Q_t|=s(P_u+P_v)$, evaluation of its diagonal kernel gives
\[
 |\rho_{Q_t}(x)|\le s\bigl(|u(x)|^2+|v(x)|^2\bigr).
\]
The density on the right has integral $2s$.  Using
\eqref{uv-orbital-potential} twice yields
\begin{equation}
 0\le D_L(\rho_{Q_t},\rho_{Q_t})\le4h_\Lambda s^2.
 \label{uv-direct}
\end{equation}
Applying the same Kato bound to the columns of $Q_t(x,y)$ gives
\begin{align}
 |\mathrm{Ex}(Q_t)|
 &\le\iint |G_L(x-y)|\|Q_t(x,y)\|_{\HS}^2\dd x\dd y\nonumber\\
 &\le\frac\pi2\Tr(|p|Q_t^2)+\frac{c_1}{L}\Tr Q_t^2
 \le2h_\Lambda s^2.
 \label{uv-exchange}
\end{align}
Consequently the pair's complete direct-minus-exchange contribution obeys
\begin{equation}
 \frac\alpha2\left[D_L(\rho_{Q_t},\rho_{Q_t})
                       -\mathrm{Ex}(Q_t)\right]
 \le3\alpha h_\Lambda s^2.
 \label{uv-pair-interaction}
\end{equation}
No sign of the exchange functional was assumed.

It remains to check the interaction with the neutralizing electron; dropping
it without an estimate would leave a gap.  The density of $P_w$ is constant,
so its direct self term and its direct cross term with $\rho_{Q_t}$ vanish.
Also $\mathrm{Ex}(P_w)=0$ by the zero-mean condition on $G_L$.  Define the
exchange bilinear form
\[
 X(Q_t,P_w)=\iint G_L(x-y)
                 \tr\bigl(Q_t(x,y)P_w(y,x)\bigr)\dd x\dd y.
\]
Writing $Q_t(p,p)$ for the diagonal momentum matrix gives exactly
\[
 X(Q_t,P_w)=\frac1V\sum_{p\ne0}\frac{4\pi}{|p|^2}
                  \langle\zeta,Q_t(p,p)\zeta\rangle.
\]
The off-diagonal terms in \eqref{uv-pair-rotation} have zero momentum
diagonal because the two Fourier supports are disjoint.  On both remaining
supports $|p|\ge3\Lambda/8$.  Normalization of $u,v$ therefore proves
\begin{equation}
 |X(Q_t,P_w)|\le\frac{512\pi}{9V\Lambda^2}s^2.
 \label{uv-spectator-exchange}
\end{equation}
The full exchange is $\mathrm{Ex}(Q_t)+2X(Q_t,P_w)$, so the extra energy is
bounded above by $\alpha|X(Q_t,P_w)|$.  This completes the accounting of
the kinetic, nuclear, direct, and exchange terms.

\subsection{Optimization and implications}
\label{uv-conclusion}

Combining \eqref{uv-kinetic}, \eqref{uv-nuclear-attraction},
\eqref{uv-pair-interaction}, and \eqref{uv-spectator-exchange} gives
\begin{align}
 \mathcal E_{0,1}(\gamma_t)
 \le m+\left[
 \left(\frac54+\frac{15\pi\alpha}{16}\right)\Lambda
       +2m+\frac{3\alpha c_1}{L}
       +\frac{512\pi\alpha}{9V\Lambda^2}
       \right]s^2
       -2\alpha b_0\Lambda sc.
 \label{uv-energy-before-optimization}
\end{align}
The last two threshold conditions in \eqref{uv-cutoff-threshold} bound the
nonlinear-in-$\Lambda$ remainder in square brackets by $\Lambda$.
Thus the square bracket is at most $A_\alpha\Lambda$.  Choose the fixed
angle
\[
 \tan t=\frac{\alpha b_0}{A_\alpha}.
\]
Then
\[
 A_\alpha s^2-2\alpha b_0sc
 =-\frac{A_\alpha\alpha^2b_0^2}
             {A_\alpha^2+\alpha^2b_0^2}
 =-c_\alpha.
\]
The variational principle and the zero Hartree--Fock vacuum energy prove
Theorem~\ref{uv-theorem}.

\paragraph{What this disproves.}
For fixed $m,\alpha,L$ and the fixed free-Wick prescription, a point-source
vacuum-polarization effect makes the optimal relative stability constant
grow at least linearly along large cutoffs:
$C_\Lambda\ge c_\alpha\Lambda-m$ already for a unit nucleus.
Combined with Theorem~\ref{thm:mf}, this also shows that the linear
cutoff order is sharp for this one-nucleus problem at fixed $\alpha,m,L$:
the energy is bounded below by $-C\alpha\Lambda$ and above by
$m-c_\alpha\Lambda$ for all sufficiently large $\Lambda$.
No subcritical condition such as $Z\alpha\le2/\pi$ repairs that statement,
since $Z=1$ suffices.  The mass cost is quadratic in the small pair rotation,
whereas its coupling to the point Coulomb source is linear; both scale with
$\Lambda$ on the chosen high-momentum packets.  All self-consistent direct
and exchange terms have been retained and bounded above, not omitted.
Keeping instead the finite periodic nuclear Robin self term changes this
example by an $O(L^{-1})$ constant at fixed $L$, and therefore does not change
the ultraviolet conclusion.

\paragraph{What this does not disprove.}
This is not a counterexample to the fixed-cutoff theorem, whose constant is
allowed to depend on $\Lambda$.  It is also not a counterexample to a
many-body bound relative to the exact interacting many-body vacuum: the
trial bounds the many-body energy from above, but does not control the
many-body vacuum subtraction, which is not the zero Hartree--Fock value.
Finally, it makes no claim about a separately constructed renormalized QED
model with cutoff-dependent bare parameters, a different reference vacuum,
or additional source renormalization.  Such choices change the Hamiltonian
or its energy convention and require their own nonperturbative analysis.
In particular no configuration-dependent subtraction has been inserted in
this proof.  The result identifies a concrete obstruction to the proposed
cutoff-uniform theorem for the model actually defined in this paper.
